% MCQ Questions for Chapter 04
\section{বহুনির্বাচনি প্রশ্ন (MCQ)}

\subsection*{\textbf{\underline{ওয়েব ডিজাইন ও HTML}}}

\begin{multicols}{2}
\begin{enumerate}

\item HTML এর পূর্ণরূপ কী? \textbf{[ঢাকা-২০২৩]}
\begin{enumerate}[label=\alph*)]
\item HyperText Markup Language
\item High Transfer Markup Language
\item HyperText Making Language
\item High Text Markup Language
\end{enumerate}

\item HTML ফাইলের এক্সটেনশন কী? \textbf{[চট্টগ্রাম-২০২৩]}
\begin{enumerate}[label=\alph*)]
\item .html বা .htm
\item .txt
\item .doc
\item .pdf
\end{enumerate}

\item HTML ট্যাগ কীভাবে লেখা হয়? \textbf{[রাজশাহী-২০২২]}
\begin{enumerate}[label=\alph*)]
\item <tag>
\item [tag]
\item (tag)
\item tag
\end{enumerate}

\item HTML ডকুমেন্টের মূল অংশ কয়টি? \textbf{[সিলেট-২০২৩]}
\begin{enumerate}[label=\alph*)]
\item ২টি (Head ও Body)
\item ৩টি
\item ৪টি
\item ৫টি
\end{enumerate}

\item <title> ট্যাগ কোথায় ব্যবহৃত হয়? \textbf{[বরিশাল-২০২২]}
\begin{enumerate}[label=\alph*)]
\item <head> এর মধ্যে
\item <body> এর মধ্যে
\item <html> এর বাইরে
\item যেকোনো জায়গায়
\end{enumerate}

\item Heading ট্যাগ কয়টি? \textbf{[দিনাজপুর-২০২৩]}
\begin{enumerate}[label=\alph*)]
\item ৬টি (h1 থেকে h6)
\item ৫টি
\item ৭টি
\item ৮টি
\end{enumerate}

\item Paragraph তৈরির ট্যাগ কোনটি? \textbf{[কুমিল্লা-২০২২]}
\begin{enumerate}[label=\alph*)]
\item <p>
\item <para>
\item <pg>
\item <paragraph>
\end{enumerate}

\item Line break এর ট্যাগ কোনটি? \textbf{[যশোর-২০২৩]}
\begin{enumerate}[label=\alph*)]
\item <br>
\item <lb>
\item <break>
\item <line>
\end{enumerate}

\end{enumerate}
\end{multicols}

\subsection*{\textbf{\underline{HTML ফরম্যাটিং ও লিস্ট}}}

\begin{multicols}{2}
\begin{enumerate}[resume]

\item Bold করার ট্যাগ কোনটি? \textbf{[ঢাকা-২০২২]}
\begin{enumerate}[label=\alph*)]
\item <b> বা <strong>
\item <bold>
\item <bd>
\item <bld>
\end{enumerate}

\item Italic করার ট্যাগ কোনটি? \textbf{[চট্টগ্রাম-২০২২]}
\begin{enumerate}[label=\alph*)]
\item <i> বা <em>
\item <italic>
\item <it>
\item <itl>
\end{enumerate}

\item Underline করার ট্যাগ কোনটি? \textbf{[রাজশাহী-২০২৩]}
\begin{enumerate}[label=\alph*)]
\item <u>
\item <underline>
\item <ul>
\item <und>
\end{enumerate}

\item Ordered List এর ট্যাগ কোনটি? \textbf{[সিলেট-২০২২]}
\begin{enumerate}[label=\alph*)]
\item <ol>
\item <ul>
\item <list>
\item <order>
\end{enumerate}

\item Unordered List এর ট্যাগ কোনটি? \textbf{[বরিশাল-২০২৩]}
\begin{enumerate}[label=\alph*)]
\item <ul>
\item <ol>
\item <list>
\item <unorder>
\end{enumerate}

\item List Item এর ট্যাগ কোনটি? \textbf{[দিনাজপুর-২০২২]}
\begin{enumerate}[label=\alph*)]
\item <li>
\item <item>
\item <list>
\item <it>
\end{enumerate}

\item Horizontal Line এর ট্যাগ কোনটি? \textbf{[কুমিল্লা-২০২৩]}
\begin{enumerate}[label=\alph*)]
\item <hr>
\item <line>
\item <hl>
\item <horizontal>
\end{enumerate}

\item Comment লেখার সিনট্যাক্স কোনটি? \textbf{[যশোর-২০২২]}
\begin{enumerate}[label=\alph*)]
\item <!-- comment -->
\item // comment
\item /* comment */
\item \# comment
\end{enumerate}

\end{enumerate}
\end{multicols}

\subsection*{\textbf{\underline{HTML লিংক ও ইমেজ}}}

\begin{multicols}{2}
\begin{enumerate}[resume]

\item Hyperlink তৈরির ট্যাগ কোনটি? \textbf{[ঢাকা-২০২৩]}
\begin{enumerate}[label=\alph*)]
\item <a>
\item <link>
\item <href>
\item <hyperlink>
\end{enumerate}

\item href এর পূর্ণরূপ কী? \textbf{[চট্টগ্রাম-২০২৩]}
\begin{enumerate}[label=\alph*)]
\item Hypertext Reference
\item High Reference
\item Hyperlink Reference
\item HTML Reference
\end{enumerate}

\item নতুন ট্যাবে লিংক খোলার attribute কোনটি? \textbf{[রাজশাহী-২০২২]}
\begin{enumerate}[label=\alph*)]
\item target="\_blank"
\item target="\_new"
\item target="\_tab"
\item target="\_open"
\end{enumerate}

\item Image যুক্ত করার ট্যাগ কোনটি? \textbf{[সিলেট-২০২৩]}
\begin{enumerate}[label=\alph*)]
\item <img>
\item <image>
\item <pic>
\item <picture>
\end{enumerate}

\item Image এর source নির্দেশ করে কোন attribute? \textbf{[বরিশাল-২০২২]}
\begin{enumerate}[label=\alph*)]
\item src
\item source
\item href
\item link
\end{enumerate}

\item Image এর alternative text এর attribute কোনটি? \textbf{[দিনাজপুর-২০২৩]}
\begin{enumerate}[label=\alph*)]
\item alt
\item text
\item title
\item name
\end{enumerate}

\item Image এর width নির্ধারণ করা হয় কীভাবে? \textbf{[কুমিল্লা-২০২২]}
\begin{enumerate}[label=\alph*)]
\item width attribute দিয়ে
\item size attribute দিয়ে
\item length attribute দিয়ে
\item dimension attribute দিয়ে
\end{enumerate}

\item Email link তৈরিতে href এ কী লিখতে হয়? \textbf{[যশোর-২০২৩]}
\begin{enumerate}[label=\alph*)]
\item mailto:email@example.com
\item email:email@example.com
\item mail:email@example.com
\item send:email@example.com
\end{enumerate}

\end{enumerate}
\end{multicols}

\subsection*{\textbf{\underline{HTML টেবিল}}}

\begin{multicols}{2}
\begin{enumerate}[resume]

\item Table তৈরির ট্যাগ কোনটি? \textbf{[ঢাকা-২০২২]}
\begin{enumerate}[label=\alph*)]
\item <table>
\item <tab>
\item <tbl>
\item <grid>
\end{enumerate}

\item Table Row এর ট্যাগ কোনটি? \textbf{[চট্টগ্রাম-২০২২]}
\begin{enumerate}[label=\alph*)]
\item <tr>
\item <row>
\item <trow>
\item <r>
\end{enumerate}

\item Table Header এর ট্যাগ কোনটি? \textbf{[রাজশাহী-২০২৩]}
\begin{enumerate}[label=\alph*)]
\item <th>
\item <thead>
\item <header>
\item <h>
\end{enumerate}

\item Table Data এর ট্যাগ কোনটি? \textbf{[সিলেট-২০২২]}
\begin{enumerate}[label=\alph*)]
\item <td>
\item <data>
\item <tdata>
\item <d>
\end{enumerate}

\item Table এ border যুক্ত করার attribute কোনটি? \textbf{[বরিশাল-২০২৩]}
\begin{enumerate}[label=\alph*)]
\item border
\item line
\item frame
\item edge
\end{enumerate}

\item Cell merge করার attribute কোনটি? \textbf{[দিনাজপুর-২০২২]}
\begin{enumerate}[label=\alph*)]
\item colspan বা rowspan
\item merge
\item combine
\item join
\end{enumerate}

\item Table caption যুক্ত করার ট্যাগ কোনটি? \textbf{[কুমিল্লা-২০২৩]}
\begin{enumerate}[label=\alph*)]
\item <caption>
\item <title>
\item <heading>
\item <cap>
\end{enumerate}

\item Table এর width নির্ধারণ করা হয় কীভাবে? \textbf{[যশোর-২০২২]}
\begin{enumerate}[label=\alph*)]
\item width attribute দিয়ে
\item size attribute দিয়ে
\item length attribute দিয়ে
\item dimension attribute দিয়ে
\end{enumerate}

\end{enumerate}
\end{multicols}

\subsection*{\textbf{\underline{HTML ফর্ম}}}

\begin{multicols}{2}
\begin{enumerate}[resume]

\item Form তৈরির ট্যাগ কোনটি? \textbf{[ঢাকা-২০২৩]}
\begin{enumerate}[label=\alph*)]
\item <form>
\item <input>
\item <field>
\item <data>
\end{enumerate}

\item Input field তৈরির ট্যাগ কোনটি? \textbf{[চট্টগ্রাম-২০২৩]}
\begin{enumerate}[label=\alph*)]
\item <input>
\item <field>
\item <text>
\item <data>
\end{enumerate}

\item Text input এর type কী? \textbf{[রাজশাহী-২০২২]}
\begin{enumerate}[label=\alph*)]
\item type="text"
\item type="input"
\item type="field"
\item type="string"
\end{enumerate}

\item Password input এর type কী? \textbf{[সিলেট-২০২৩]}
\begin{enumerate}[label=\alph*)]
\item type="password"
\item type="pass"
\item type="pwd"
\item type="secret"
\end{enumerate}

\item Radio button এর type কী? \textbf{[বরিশাল-২০২২]}
\begin{enumerate}[label=\alph*)]
\item type="radio"
\item type="button"
\item type="option"
\item type="select"
\end{enumerate}

\item Checkbox এর type কী? \textbf{[দিনাজপুর-২০২৩]}
\begin{enumerate}[label=\alph*)]
\item type="checkbox"
\item type="check"
\item type="box"
\item type="tick"
\end{enumerate}

\item Submit button এর type কী? \textbf{[কুমিল্লা-২০২২]}
\begin{enumerate}[label=\alph*)]
\item type="submit"
\item type="send"
\item type="button"
\item type="ok"
\end{enumerate}

\item Dropdown list তৈরির ট্যাগ কোনটি? \textbf{[যশোর-২০২৩]}
\begin{enumerate}[label=\alph*)]
\item <select>
\item <dropdown>
\item <list>
\item <option>
\end{enumerate}

\end{enumerate}
\end{multicols}

\subsection*{\textbf{\underline{CSS বেসিক}}}

\begin{multicols}{2}
\begin{enumerate}[resume]

\item CSS এর পূর্ণরূপ কী? \textbf{[ঢাকা-২০২২]}
\begin{enumerate}[label=\alph*)]
\item Cascading Style Sheets
\item Computer Style Sheets
\item Creative Style Sheets
\item Colorful Style Sheets
\end{enumerate}

\item Inline CSS কীভাবে লেখা হয়? \textbf{[চট্টগ্রাম-২০২২]}
\begin{enumerate}[label=\alph*)]
\item style attribute দিয়ে
\item <style> ট্যাগে
\item আলাদা ফাইলে
\item <css> ট্যাগে
\end{enumerate}

\item Internal CSS কোথায় লেখা হয়? \textbf{[রাজশাহী-২০২৩]}
\begin{enumerate}[label=\alph*)]
\item <head> এর মধ্যে <style> ট্যাগে
\item <body> এর মধ্যে
\item আলাদা ফাইলে
\item যেকোনো জায়গায়
\end{enumerate}

\item External CSS ফাইলের এক্সটেনশন কী? \textbf{[সিলেট-২০২২]}
\begin{enumerate}[label=\alph*)]
\item .css
\item .style
\item .cs
\item .stylesheet
\end{enumerate}

\item CSS এ color নির্ধারণ করা যায় কীভাবে? \textbf{[বরিশাল-২০২৩]}
\begin{enumerate}[label=\alph*)]
\item color: red; বা color: \#FF0000;
\item color = red;
\item color(red);
\item red color;
\end{enumerate}

\item Background color নির্ধারণের property কোনটি? \textbf{[দিনাজপুর-২০২২]}
\begin{enumerate}[label=\alph*)]
\item background-color
\item bg-color
\item color-background
\item bgcolor
\end{enumerate}

\item Font size নির্ধারণের property কোনটি? \textbf{[কুমিল্লা-২০২৩]}
\begin{enumerate}[label=\alph*)]
\item font-size
\item size
\item text-size
\item font
\end{enumerate}

\item Text align করার property কোনটি? \textbf{[যশোর-২০২২]}
\begin{enumerate}[label=\alph*)]
\item text-align
\item align
\item alignment
\item text-position
\end{enumerate}

\end{enumerate}
\end{multicols}

\subsection*{\textbf{\underline{উত্তরমালা}}}

\begin{multicols}{4}
\begin{enumerate}
\item a
\item a
\item a
\item a
\item a
\item a
\item a
\item a
\item a
\item a
\item a
\item a
\item a
\item a
\item a
\item a
\item a
\item a
\item a
\item a
\item a
\item a
\item a
\item a
\item a
\item a
\item a
\item a
\item a
\item a
\item a
\item a
\item a
\item a
\item a
\item a
\item a
\item a
\item a
\item a
\item a
\item a
\item a
\item a
\item a
\item a
\item a
\item a
\item a
\item a
\item a
\item a
\item a
\item a
\item a
\item a
\end{enumerate}
\end{multicols}

\subsection*{\textbf{\underline{CSS Advanced}}}

\begin{multicols}{2}
\begin{enumerate}[resume]

\item Class selector কীভাবে লেখা হয়? \textbf{[ঢাকা-২০২৩]}
\begin{enumerate}[label=\alph*)]
\item .classname
\item \#classname
\item classname
\item *classname
\end{enumerate}

\item ID selector কীভাবে লেখা হয়? \textbf{[চট্টগ্রাম-২০২২]}
\begin{enumerate}[label=\alph*)]
\item \#idname
\item .idname
\item idname
\item *idname
\end{enumerate}

\item Universal selector কোনটি? \textbf{[সিলেট-২০২৩]}
\begin{enumerate}[label=\alph*)]
\item *
\item \#
\item .
\item @
\end{enumerate}

\item Margin property কী কাজ করে? \textbf{[রাজশাহী-২০২২]}
\begin{enumerate}[label=\alph*)]
\item বাইরের স্পেস নিয়ন্ত্রণ
\item ভিতরের স্পেস নিয়ন্ত্রণ
\item বর্ডার নিয়ন্ত্রণ
\item রং নিয়ন্ত্রণ
\end{enumerate}

\item Padding property কী কাজ করে? \textbf{[বরিশাল-২০২৩]}
\begin{enumerate}[label=\alph*)]
\item ভিতরের স্পেস নিয়ন্ত্রণ
\item বাইরের স্পেস নিয়ন্ত্রণ
\item বর্ডার নিয়ন্ত্রণ
\item রং নিয়ন্ত্রণ
\end{enumerate}

\item Border property কী কাজ করে? \textbf{[দিনাজপুর-২০২২]}
\begin{enumerate}[label=\alph*)]
\item সীমানা তৈরি
\item স্পেস তৈরি
\item রং পরিবর্তন
\item আকার পরিবর্তন
\end{enumerate}

\item Display property এর মান কয়টি? \textbf{[কুমিল্লা-২০২৩]}
\begin{enumerate}[label=\alph*)]
\item অনেকগুলো (block, inline, none ইত্যাদি)
\item ২টি
\item ৩টি
\item ৪টি
\end{enumerate}

\item Position property এর মান কয়টি? \textbf{[যশোর-২০২২]}
\begin{enumerate}[label=\alph*)]
\item ৫টি (static, relative, absolute, fixed, sticky)
\item ২টি
\item ৩টি
\item ৪টি
\end{enumerate}

\item Float property কী কাজ করে? \textbf{[ঢাকা-২০২৩]}
\begin{enumerate}[label=\alph*)]
\item Element কে left/right এ সরানো
\item Element কে up/down এ সরানো
\item Element লুকানো
\item Element দেখানো
\end{enumerate}

\item Z-index কী কাজ করে? \textbf{[চট্টগ্রাম-২০২২]}
\begin{enumerate}[label=\alph*)]
\item Element এর স্তর নিয়ন্ত্রণ
\item Element এর রং নিয়ন্ত্রণ
\item Element এর আকার নিয়ন্ত্রণ
\item Element এর অবস্থান নিয়ন্ত্রণ
\end{enumerate}

\item Overflow property কী কাজ করে? \textbf{[সিলেট-২০২৩]}
\begin{enumerate}[label=\alph*)]
\item অতিরিক্ত content নিয়ন্ত্রণ
\item রং নিয়ন্ত্রণ
\item আকার নিয়ন্ত্রণ
\item অবস্থান নিয়ন্ত্রণ
\end{enumerate}

\item Opacity property এর মান কত থেকে কত? \textbf{[রাজশাহী-২০২২]}
\begin{enumerate}[label=\alph*)]
\item 0 থেকে 1
\item 0 থেকে 10
\item 0 থেকে 100
\item 1 থেকে 10
\end{enumerate}

\item Box-shadow কী কাজ করে? \textbf{[বরিশাল-২০২৩]}
\begin{enumerate}[label=\alph*)]
\item ছায়া তৈরি
\item রং পরিবর্তন
\item আকার পরিবর্তন
\item অবস্থান পরিবর্তন
\end{enumerate}

\item Text-shadow কী কাজ করে? \textbf{[দিনাজপুর-২০২২]}
\begin{enumerate}[label=\alph*)]
\item টেক্সটে ছায়া
\item টেক্সট রং
\item টেক্সট আকার
\item টেক্সট অবস্থান
\end{enumerate}

\item Border-radius কী কাজ করে? \textbf{[কুমিল্লা-২০২৩]}
\begin{enumerate}[label=\alph*)]
\item কোণা গোলাকার করা
\item বর্ডার রং
\item বর্ডার আকার
\item বর্ডার অবস্থান
\end{enumerate}

\item Gradient কী? \textbf{[যশোর-২০২২]}
\begin{enumerate}[label=\alph*)]
\item রঙের মিশ্রণ
\item একটি রং
\item কোনো রং নয়
\item সাদা রং
\end{enumerate}

\end{enumerate}
\end{multicols}

\subsection*{\textbf{\underline{JavaScript বেসিক}}}

\begin{multicols}{2}
\begin{enumerate}[resume]

\item JavaScript কী ধরনের ভাষা? \textbf{[ঢাকা, চট্টগ্রাম-২০২৩]}
\begin{enumerate}[label=\alph*)]
\item Scripting Language
\item Markup Language
\item Style Language
\item Database Language
\end{enumerate}

\item JavaScript কোথায় লেখা হয়? \textbf{[সিলেট-২০২২]}
\begin{enumerate}[label=\alph*)]
\item <script> ট্যাগে
\item <style> ট্যাগে
\item <code> ট্যাগে
\item <js> ট্যাগে
\end{enumerate}

\item Variable declare করার keyword কোনটি? \textbf{[রাজশাহী-২০২৩]}
\begin{enumerate}[label=\alph*)]
\item var, let, const
\item variable
\item declare
\item define
\end{enumerate}

\item Comment লেখার সিনট্যাক্স কোনটি? \textbf{[বরিশাল-২০২২]}
\begin{enumerate}[label=\alph*)]
\item // comment বা /* comment */
\item <!-- comment -->
\item \# comment
\item ' comment
\end{enumerate}

\item Alert box দেখানোর function কোনটি? \textbf{[দিনাজপুর-২০২৩]}
\begin{enumerate}[label=\alph*)]
\item alert()
\item message()
\item show()
\item display()
\end{enumerate}

\item Console এ লেখার function কোনটি? \textbf{[কুমিল্লা-২০২২]}
\begin{enumerate}[label=\alph*)]
\item console.log()
\item print()
\item write()
\item display()
\end{enumerate}

\item Prompt box দেখানোর function কোনটি? \textbf{[যশোর-২০২৩]}
\begin{enumerate}[label=\alph*)]
\item prompt()
\item input()
\item ask()
\item get()
\end{enumerate}

\item Confirm box দেখানোর function কোনটি? \textbf{[ঢাকা-২০২২]}
\begin{enumerate}[label=\alph*)]
\item confirm()
\item check()
\item verify()
\item validate()
\end{enumerate}

\item String concatenation এর operator কোনটি? \textbf{[চট্টগ্রাম-২০২২]}
\begin{enumerate}[label=\alph*)]
\item +
\item \&
\item *
\item .
\end{enumerate}

\item Equality check এর operator কোনটি? \textbf{[সিলেট-২০২৩]}
\begin{enumerate}[label=\alph*)]
\item == বা ===
\item =
\item !=
\item <>
\end{enumerate}

\item Logical AND operator কোনটি? \textbf{[রাজশাহী-২০২২]}
\begin{enumerate}[label=\alph*)]
\item \&\&
\item ||
\item !
\item \&
\end{enumerate}

\item Logical OR operator কোনটি? \textbf{[বরিশাল-২০২৩]}
\begin{enumerate}[label=\alph*)]
\item ||
\item \&\&
\item !
\item |
\end{enumerate}

\item Logical NOT operator কোনটি? \textbf{[দিনাজপুর-২০২২]}
\begin{enumerate}[label=\alph*)]
\item !
\item \&\&
\item ||
\item ~
\end{enumerate}

\item If statement এর সিনট্যাক্স কোনটি? \textbf{[কুমিল্লা-২০২৩]}
\begin{enumerate}[label=\alph*)]
\item if (condition) \{ \}
\item if condition then
\item if: condition
\item if [condition]
\end{enumerate}

\item For loop এর সিনট্যাক্স কোনটি? \textbf{[যশোর-২০২২]}
\begin{enumerate}[label=\alph*)]
\item for (i=0; i<10; i++) \{ \}
\item for i in range(10)
\item for i to 10
\item for (i:10)
\end{enumerate}

\item While loop এর সিনট্যাক্স কোনটি? \textbf{[ঢাকা-২০২৩]}
\begin{enumerate}[label=\alph*)]
\item while (condition) \{ \}
\item while condition do
\item while: condition
\item while [condition]
\end{enumerate}

\item Function declare করার সিনট্যাক্স কোনটি? \textbf{[চট্টগ্রাম-২০২২]}
\begin{enumerate}[label=\alph*)]
\item function name() \{ \}
\item def name():
\item function: name
\item func name()
\end{enumerate}

\item Array declare করার সিনট্যাক্স কোনটি? \textbf{[সিলেট-২০২৩]}
\begin{enumerate}[label=\alph*)]
\item var arr = []
\item array arr
\item arr = array()
\item arr[]
\end{enumerate}

\item Object declare করার সিনট্যাক্স কোনটি? \textbf{[রাজশাহী-২০২২]}
\begin{enumerate}[label=\alph*)]
\item var obj = \{ \}
\item object obj
\item obj = object()
\item obj\{ \}
\end{enumerate}

\item DOM এর পূর্ণরূপ কী? \textbf{[বরিশাল-২০২৩]}
\begin{enumerate}[label=\alph*)]
\item Document Object Model
\item Data Object Model
\item Digital Object Model
\item Dynamic Object Model
\end{enumerate}

\item getElementById() কী করে? \textbf{[দিনাজপুর-২০২২]}
\begin{enumerate}[label=\alph*)]
\item ID দিয়ে element খুঁজে
\item Class দিয়ে element খুঁজে
\item Tag দিয়ে element খুঁজে
\item Name দিয়ে element খুঁজে
\end{enumerate}

\item getElementsByClassName() কী করে? \textbf{[কুমিল্লা-২০২৩]}
\begin{enumerate}[label=\alph*)]
\item Class দিয়ে element খুঁজে
\item ID দিয়ে element খুঁজে
\item Tag দিয়ে element খুঁজে
\item Name দিয়ে element খুঁজে
\end{enumerate}

\item Event কী? \textbf{[যশোর-২০২২]}
\begin{enumerate}[label=\alph*)]
\item ব্যবহারকারীর কার্যক্রম
\item একটি function
\item একটি variable
\item একটি object
\end{enumerate}

\item onclick event কখন ঘটে? \textbf{[ঢাকা-২০২৩]}
\begin{enumerate}[label=\alph*)]
\item ক্লিক করলে
\item মাউস রাখলে
\item টাইপ করলে
\item স্ক্রল করলে
\end{enumerate}

\item onload event কখন ঘটে? \textbf{[চট্টগ্রাম-২০২২]}
\begin{enumerate}[label=\alph*)]
\item পেজ লোড হলে
\item ক্লিক করলে
\item টাইপ করলে
\item স্ক্রল করলে
\end{enumerate}

\end{enumerate}
\end{multicols}

\subsection*{\textbf{\underline{উত্তরমালা (৫৭-১০৬)}}}

\begin{multicols}{4}
\begin{enumerate}[start=57]
\item a
\item a
\item a
\item a
\item a
\item a
\item a
\item a
\item a
\item a
\item a
\item a
\item a
\item a
\item a
\item a
\item a
\item a
\item a
\item a
\item a
\item a
\item a
\item a
\item a
\item a
\item a
\item a
\item a
\item a
\item a
\item a
\item a
\item a
\item a
\item a
\item a
\item a
\item a
\item a
\item a
\item a
\item a
\item a
\item a
\item a
\item a
\item a
\item a
\item a
\end{enumerate}
\end{multicols}
